%% 自動生成（tools/mkchecklist.py）。各章の「ふりかえり」の項目を集めたもの
\begingroup\footnotesize
\def\bkcb{\tikz[baseline=-0.5ex]\draw[\bkTF, line width=0.7pt, rounded corners=0.5pt] (0,-0.3em) rectangle (0.8em,0.5em);}
\begin{longtblr}{colspec={Q[c,m,wd=9mm] X[l,m] Q[c,m,wd=10mm] Q[c,m,wd=9mm] Q[c,m,wd=9mm]}, width=\linewidth, rowhead=1, colsep=3pt,
  row{1}={bg=bkVol!14, font=\sffamily\bfseries\scriptsize}, hlines={0.4pt,bkLine}, rowsep=0.7pt}
章 & できるようになったこと & ページ & 1回め & 2回め \\
{\sffamily\bfseries\color{\bkTT}6-1} & 微分方程式を、ことばと単位で読み、式に書ける。 & \pageref{bk:6-1:yomu} & \bkcb & \bkcb \\
 & 傾きの場から、解のようすを予想できる。 & \pageref{bk:6-1:ba} & \bkcb & \bkcb \\
 & 代入して、解かどうか確かめられる。 & \pageref{bk:6-1:tashikame} & \bkcb & \bkcb \\
{\sffamily\bfseries\color{\bkTT}6-2} & 「次の値＝今の値＋今の変わる速さ×$\Delta t$」で、表計算で解ける。 & \pageref{bk:6-2:ippo} & \bkcb & \bkcb \\
 & $\Delta t$ を小さくすると正確になることを確かめられる。 & \pageref{bk:6-2:zure} & \bkcb & \bkcb \\
 & ばねの式を、位置と速度の2つの列で、正しい順に解ける。 & \pageref{bk:6-2:bane} & \bkcb & \bkcb \\
{\sffamily\bfseries\color{\bkTT}6-3} & 右辺を見て、積分するだけで解ける形を見分け、解ける。 & \pageref{bk:6-3:miwake} & \bkcb & \bkcb \\
 & 深さや場所が変数のルールも、積分して解ける。 & \pageref{bk:6-3:basho} & \bkcb & \bkcb \\
 & 両はしの条件（境界条件）から、積分定数を決められる。 & \pageref{bk:6-3:kyokai} & \bkcb & \bkcb \\
{\sffamily\bfseries\color{\bkTT}6-4} & 変数分離で、自分に比例してへる式を解ける。 & \pageref{bk:6-4:bunri} & \bkcb & \bkcb \\
 & 時定数・半減期を使って、かかる時間を求められる。 & \pageref{bk:6-4:tau} & \bkcb & \bkcb \\
 & RC・RL の回路の式を立てて、時定数を求められる。 & \pageref{bk:6-4:rc} & \bkcb & \bkcb \\
{\sffamily\bfseries\color{\bkTT}6-5} & 落ちる物の式を立て、終端速度を求められる。 & \pageref{bk:6-5:shuutan} & \bkcb & \bkcb \\
 & 表計算で、落ちる物の速さを計算できる。 & \pageref{bk:6-5:hyou} & \bkcb & \bkcb \\
 & 速さに比例する抵抗の式を解ける（RL 回路も）。 & \pageref{bk:6-5:toku} & \bkcb & \bkcb \\
{\sffamily\bfseries\color{\bkTT}6-6} & ばねの運動方程式を立て、$\omega_{\mathrm n}=\sqrt{k/m}$ と周期を求められる。 & \pageref{bk:6-6:shiki} & \bkcb & \bkcb \\
 & はじめの位置と速度から、振幅と初期位相を決められる。 & \pageref{bk:6-6:shoki} & \bkcb & \bkcb \\
 & 振り子・LC回路の周期を、ばねと同じ形の式から求められる。 & \pageref{bk:6-6:lc} & \bkcb & \bkcb \\
{\sffamily\bfseries\color{\bkTT}6-7} & ばね・質量・ダンパーの式を、力の図から立てられる。 & \pageref{bk:6-7:shiki} & \bkcb & \bkcb \\
 & となりの山の比から、対数減衰率を求められる。 & \pageref{bk:6-7:yama} & \bkcb & \bkcb \\
 & 減衰比 $\zeta$ を計算して、止まり方を答えられる（RLC回路でも）。 & \pageref{bk:6-7:zeta} & \bkcb & \bkcb \\
{\sffamily\bfseries\color{\bkTT}6-8} & $j$ をかけると90°回ることから、$j^2=-1$ を説明できる。 & \pageref{bk:6-8:kaiten} & \bkcb & \bkcb \\
 & 複素数のたし算・ひき算・かけ算・わり算ができる。 & \pageref{bk:6-8:kakeru} & \bkcb & \bkcb \\
 & 判別式がマイナスの2次方程式を解いて、虚数の解を求められる。 & \pageref{bk:6-8:kai} & \bkcb & \bkcb \\
{\sffamily\bfseries\color{\bkTT}6-9} & 複素数を平面の点で表し、絶対値と偏角を求められる。 & \pageref{bk:6-9:kyoku} & \bkcb & \bkcb \\
 & かけ算が「大きさはかけ算、角はたし算」になることを説明できる。 & \pageref{bk:6-9:kakeru} & \bkcb & \bkcb \\
 & かけ算で、点を回したり、うでの先の位置を求めたりできる。 & \pageref{bk:6-9:kaiten} & \bkcb & \bkcb \\
{\sffamily\bfseries\color{\bkTT}6-10} & $e^{j\theta}$ が単位円の上の角 $\theta$ の点になるわけを、小さな回転のくり返しと微分で説明できる。 & \pageref{bk:6-10:kaiten} & \bkcb & \bkcb \\
 & $e^{j\theta}$ を $a+bj$ に直し、$re^{j\theta}$ のかけ算・わり算・べきを計算できる。 & \pageref{bk:6-10:kyoku} & \bkcb & \bkcb \\
 & フェーザのたし算で交流を合成し、微分が $j\omega$ をかけることだと説明できる。 & \pageref{bk:6-10:phasor} & \bkcb & \bkcb \\
{\sffamily\bfseries\color{\bkTT}6-11} & $\cosh x$・$\sinh x$・$\tanh x$ を $e^x$ と $e^{-x}$ で表し、値とグラフの形を説明できる。 & \pageref{bk:6-11:teigi} & \bkcb & \bkcb \\
 & $\cosh^2x-\sinh^2x=1$ と、微分のきまりを使える。 & \pageref{bk:6-11:bibun} & \bkcb & \bkcb \\
 & $y''=a^2y$ の答えを $\cosh$・$\sinh$ で書き、落下や電線の形に使える。 & \pageref{bk:6-11:kaku} & \bkcb & \bkcb \\
{\sffamily\bfseries\color{\bkTT}6-12} & 式に $e^{st}$ を入れて、特性方程式を作れる。 & \pageref{bk:6-12:oku} & \bkcb & \bkcb \\
 & 解が2つ・複素数・重解のとき、答えを組み立てられる。 & \pageref{bk:6-12:ba3} & \bkcb & \bkcb \\
 & $\zeta$ で動き方を見分け、車や RLC 回路の数を求められる。 & \pageref{bk:6-12:zeta} & \bkcb & \bkcb \\
{\sffamily\bfseries\color{\bkTT}6-13} & ゆすり続けたときのゆれを、表計算と共振曲線で説明できる。 & \pageref{bk:6-13:kyoshin} & \bkcb & \bkcb \\
 & ゆすりに合うゆれを、係数を比べる方法と複素数で求められる。 & \pageref{bk:6-13:keisu} & \bkcb & \bkcb \\
 & いちばん大きくゆれる $\omega$ を微分で求め、RLC 回路の共振周波数を求められる。 & \pageref{bk:6-13:yama} & \bkcb & \bkcb \\
{\sffamily\bfseries\color{\bkTT}6-14} & からみあう2つの量を、表計算で同時に進めて計算できる。 & \pageref{bk:6-14:hyou} & \bkcb & \bkcb \\
 & 和と差（重みつきの和と差）に分けて、2つの式を解ける。 & \pageref{bk:6-14:wasa} & \bkcb & \bkcb \\
 & 振動モードや消去法で、連立微分方程式の答えを求められる。 & \pageref{bk:6-14:mode} & \bkcb & \bkcb \\
{\sffamily\bfseries\color{\bkTT}6-15} & 法則から、微分方程式とはじめの状態を書ける。 & \pageref{bk:6-15:buhin} & \bkcb & \bkcb \\
 & 相似の表で、ばね・電気・熱の量を対応させられる。 & \pageref{bk:6-15:soji} & \bkcb & \bkcb \\
 & 式の形から、解き方（章）を選べる。 & \pageref{bk:6-15:chizu} & \bkcb & \bkcb \\
\end{longtblr}
\endgroup
